\chapter*{Preface}
\addcontentsline{toc}{chapter}{Preface}
\markboth{Preface}{Preface}

Programming is a conversation with a computer. You want something done, you
write it down in a language the computer understands, and the computer does
exactly that, without argument and at an unbelievable speed. Learning to hold
that conversation is easier than most people think, as long as the first
steps are small and every step is clear. That is what this book sets out to
give you.

The language we use is Salam. It is a real, complete programming language,
not an educational toy: large programs are written in it, web pages and sites
are built with it, and the Salam compiler itself is written in Salam. Its
words are short and plain, its rules are strict in ways that catch mistakes
early, and its error messages say not only what went wrong but how to fix
it. Salam is also bilingual: a program can be written in English or in
Persian, with exactly the same power and the same speed. This book uses the
English spelling throughout. A Persian edition of the book covers the same
ground in the same order.

\section*{Who is this book for?}

For anyone who wants to learn programming from nothing. There are no
prerequisites: no earlier experience with programming, no advanced
mathematics. All you need is a computer, or even a phone with a web browser,
and a little patience. Secondary school students, first-year university
students, teachers, and every curious adult who has always wanted to know
what programming is about are the readers we had in mind. If you have
programmed in another language before, you can read this book too, but you
will probably move through the first chapters quickly.

\section*{Install nothing}

One of our firm decisions in writing this book was that the reader should
never have to deal with a command line, a terminal or a software
installation. You can write and run every program in the book in the Salam
online editor:
\begin{center}
\url{https://editor.salamlang.ir}
\end{center}
The editor runs in your browser, depends on no other service, and is the
same tool you will meet in the first chapter. Open the address, choose
English as the language, and start writing.

\section*{What will you learn?}

The path of this book is the familiar path of every good programming book,
with one difference: at every step you write real Salam code.
\begin{itemize}
  \item In chapter \ref{ch:first} you write your first program and learn
        how to put something on the screen.
  \item Chapters \ref{ch:variables} to \ref{ch:operators} cover variables,
        memory, data types and arithmetic: the building materials of every
        program.
  \item In chapter \ref{ch:conditions} programs learn to make decisions, and
        in chapter \ref{ch:loops} they learn to do something many times
        over.
  \item Chapter \ref{ch:functions} shows you how to split a program into
        small, reusable pieces.
  \item Chapters \ref{ch:arrays} and \ref{ch:strings} are about keeping
        lists of data and working with text.
  \item Chapter \ref{ch:projects} brings all of this together in a few
        small, practical programs.
\end{itemize}
We have deliberately stayed away from famous algorithms and long programs.
The aim of this book is for you to learn the basics so well that afterwards
you can learn anything else on your own. The examples are all simple and
concrete: a grade average, a shopping bill, a temperature table, a
countdown, and things of that kind.

\section*{How to read this book}

Programming is not learned by reading. It is learned by writing. We ask you
to type every program you see in this book into the editor yourself (do not
copy and paste it), run it, and compare the output with what the book shows.
Then tamper with it: change a number, delete a line, add something, and see
what happens. Experimenting is the best teacher you will ever have.

Throughout the book you will see several kinds of boxes. \textbf{Note}
boxes make something clearer, \textbf{Tip} boxes suggest a better way of
doing things, \textbf{Watch out} boxes warn you about common mistakes, and
\textbf{Going further} boxes are written for curious readers and can be
skipped on a first reading. Every chapter ends with exercises. We have
included the answers to a selection of them in the last appendix, but please
make an honest attempt before you look.

\section*{A word from the authors}

This book is the result of a collaboration between the creator of the Salam
language and a university lecturer. One of us has looked at the language
from the inside, the other from the point of view of the classroom and of a
student meeting programming for the first time. We hope that these two views,
side by side, have produced a book that is both correct and kind. If you find
a mistake while reading, or have a suggestion, we would be glad to hear from
you through the Salam project page on GitHub.

\begin{flushright}
Seyyed Ali Mohammadiyeh\\
Mahdi Esmaeili
\end{flushright}
\clearpage
